\chapter{Simple Harmonic Motion (Dynamics)}

\section{Introduction and Governing Differential Equation}

\begin{defbox}[Simple Harmonic Motion (SHM)]
A particle moving along a straight line is said to execute \textbf{Simple Harmonic Motion} if its acceleration is always directed towards a fixed point on the line (called the \textbf{center of oscillation} or \textbf{mean position}) and its magnitude is directly proportional to the displacement of the particle from that fixed point.
\end{defbox}

Let $O$ be the mean position on line $x' O x$, and let $P$ be the position of the particle at time $t$ with displacement $OP = x$.
The governing differential equation of motion per unit mass is:
\begin{equation}
    \frac{d^2 x}{d t^2} = -\mu x \iff \ddot{x} + \mu x = 0
\end{equation}
where $\mu > 0$ is the constant of proportionality representing the acceleration per unit displacement.

\begin{figure}[htbp]
\centering
\begin{tikzpicture}[scale=1.4, >=Stealth]
    % Line of motion
    \draw[thick, darkslate] (-3.5, 0) -- (3.5, 0);
    
    % Main Points on Axis
    \filldraw[accentred] (-2.5, 0) circle (2.5pt) node[below=4pt] {$A'$ (Extreme)};
    \filldraw[primaryblue] (0, 0) circle (2.5pt) node[below=4pt] {$O$ (Mean Center)};
    \filldraw[accentred] (2.5, 0) circle (2.5pt) node[below=4pt] {$A$ (Extreme)};
    \filldraw[accentgreen] (1.5, 0) circle (2.5pt) node[above right=2pt] {$P$};
    
    % Dashed vertical projection lines
    \draw[dashed, bordergrey] (0, 0) -- (0, 1.4);
    \draw[dashed, bordergrey] (1.5, 0) -- (1.5, 1.4);
    \draw[dashed, bordergrey] (2.5, 0) -- (2.5, -1.0);
    \draw[dashed, bordergrey] (0, 0) -- (0, -1.0);
    
    % Displacement x Arrow
    \draw[->, thick, primaryblue] (0, 0.6) -- (1.5, 0.6) node[midway, above] {Displacement $x$};
    
    % Acceleration vector pointing towards O
    \draw[->, ultra thick, accentred] (1.5, 1.2) -- (0.3, 1.2) node[midway, above] {Acceleration $\mathbf{a}_x = -\mu x$};
    
    % Amplitude a Dimension Line
    \draw[<->, thick, darkslate] (0, -0.9) -- (2.5, -0.9) node[midway, fill=white, inner sep=2pt] {Amplitude $a$};
\end{tikzpicture}
\caption{Kinematics of a particle executing Simple Harmonic Motion about mean position $O$.}
\end{figure}

\section{Kinematic Expressions}

\subsection{Velocity as a Function of Position}
Multiplying $\ddot{x} = -\mu x$ by $2 \dot{x}$ and integrating with respect to $t$:
\begin{equation}
    2 \dot{x} \ddot{x} = -2\mu x \dot{x} \implies \int \frac{d}{dt}(\dot{x}^2) \, dt = -2\mu \int x \, dx \implies v^2 = -\mu x^2 + C
\end{equation}
Let $a$ be the maximum displacement (\textbf{amplitude}) where the particle comes to instantaneous rest ($v = 0$ at $x = a$). Then $0 = -\mu a^2 + C \implies C = \mu a^2$.
\begin{equation}
    v^2 = \mu (a^2 - x^2) \implies v = \pm \sqrt{\mu} \sqrt{a^2 - x^2}
\end{equation}

\subsection{Displacement as a Function of Time}
Choosing the positive sign for motion away from origin $O$:
\begin{equation}
    \frac{dx}{dt} = \sqrt{\mu}\sqrt{a^2 - x^2} \implies \int \frac{dx}{\sqrt{a^2 - x^2}} = \int \sqrt{\mu} \, dt
\end{equation}
\begin{itemize}
    \item If time $t=0$ is measured when particle is at extreme position $A$ ($x = a$):
    \begin{equation}
        x(t) = a \cos(\sqrt{\mu} t)
    \end{equation}
    \item If time $t=0$ is measured when particle passes through origin $O$ ($x = 0$):
    \begin{equation}
        x(t) = a \sin(\sqrt{\mu} t)
    \end{equation}
\end{itemize}

\section{Important Properties of SHM}

\begin{keybox}[Key Physical Quantities in SHM]
\begin{enumerate}[leftmargin=1.5em]
    \item \textbf{Amplitude ($a$)}: Maximum displacement from the mean center $O$.
    \item \textbf{Time Period ($T$)}: The time taken for one complete oscillation:
    \begin{equation}
        T = \frac{2\pi}{\sqrt{\mu}}
    \end{equation}
    \item \textbf{Frequency ($n$)}: Number of complete oscillations per unit time:
    \begin{equation}
        n = \frac{1}{T} = \frac{\sqrt{\mu}}{2\pi}
    \end{equation}
    \item \textbf{Maximum Speed}: $v_{\max} = a \sqrt{\mu} = \dfrac{2\pi a}{T}$ (achieved at mean position $x=0$).
    \item \textbf{Maximum Acceleration}: $f_{\max} = a \mu = \dfrac{4\pi^2 a}{T^2}$ (achieved at extreme positions $x = \pm a$).
\end{enumerate}
\end{keybox}

\section{Worked Examples (Complete Unabridged Collection)}

\begin{examplebox}{Period from Velocities at Two Positions}
A point moving in a straight line with SHM has velocities $v_1$ and $v_2$ when its distances from the center are $x_1$ and $x_2$ respectively. Show that the period of motion is:
\begin{equation*}
    T = 2\pi \sqrt{\frac{x_1^2 - x_2^2}{v_2^2 - v_1^2}}
\end{equation*}

\tcblower
\textbf{Solution:}

\begin{enumerate}[leftmargin=1.5em]
    \item \textbf{Apply Velocity-Displacement Relation}:\\
    For SHM, the relation between velocity $v$ and position $x$ is:
    \begin{align}
        v_1^2 &= \mu (a^2 - x_1^2) \label{eq:ex1_v1} \\
        v_2^2 &= \mu (a^2 - x_2^2) \label{eq:ex1_v2}
    \end{align}
    
    \item \textbf{Eliminate Amplitude $a$}:\\
    Subtracting equation \eqref{eq:ex1_v1} from equation \eqref{eq:ex1_v2}:
    \begin{equation*}
        v_2^2 - v_1^2 = \mu(a^2 - x_2^2) - \mu(a^2 - x_1^2) = \mu(x_1^2 - x_2^2)
    \end{equation*}
    Solving for intensity constant $\mu$:
    \begin{equation*}
        \mu = \frac{v_2^2 - v_1^2}{x_1^2 - x_2^2} \implies \sqrt{\mu} = \sqrt{\frac{v_2^2 - v_1^2}{x_1^2 - x_2^2}}
    \end{equation*}
    
    \item \textbf{Calculate Time Period $T$}:\\
    Using the standard period formula $T = \dfrac{2\pi}{\sqrt{\mu}}$:
    \begin{equation*}
        T = \frac{2\pi}{\sqrt{\frac{v_2^2 - v_1^2}{x_1^2 - x_2^2}}} = 2\pi \sqrt{\frac{x_1^2 - x_2^2}{v_2^2 - v_1^2}}
    \end{equation*}
\end{enumerate}
\end{examplebox}

\begin{examplebox}{State Parameters Calculation}
A particle executing SHM has velocities $u$ and $v$ at two positions where its accelerations are $\alpha$ and $\beta$ respectively. Show that the distance between these two positions is $\dfrac{v^2 - u^2}{\alpha + \beta}$ and the amplitude of motion is:
\begin{equation*}
    a = \sqrt{\frac{(v^2 - u^2)(\alpha^2 v^2 - \beta^2 u^2)}{(\beta^2 - \alpha^2)^2}}
\end{equation*}

\tcblower
\textbf{Solution:}

\begin{enumerate}[leftmargin=1.5em]
    \item \textbf{Express Accelerations and Positions}:\\
    Let the displacements corresponding to velocities $u$ and $v$ be $x_1$ and $x_2$.
    The magnitude of acceleration at position $x$ is $f = \mu x$:
    \begin{align*}
        \alpha &= \mu x_1 \implies x_1 = \frac{\alpha}{\mu} \\
        \beta &= \mu x_2 \implies x_2 = \frac{\beta}{\mu}
    \end{align*}
    
    \item \textbf{Determine Distance Between Positions}:\\
    The distance between the two points is $d = x_2 - x_1$:
    \begin{equation*}
        d = x_2 - x_1 = \frac{\beta - \alpha}{\mu}
    \end{equation*}
    Using $u^2 = \mu a^2 - \mu x_1^2 = \mu a^2 - \frac{\alpha^2}{\mu}$ and $v^2 = \mu a^2 - \frac{\beta^2}{\mu}$:
    \begin{equation*}
        v^2 - u^2 = \frac{\alpha^2 - \beta^2}{\mu} = \frac{(\alpha - \beta)(\alpha + \beta)}{\mu}
    \end{equation*}
    Dividing by $(\alpha + \beta)$:
    \begin{equation*}
        \frac{v^2 - u^2}{\alpha + \beta} = \frac{\alpha - \beta}{\mu} = x_2 - x_1 = d
    \end{equation*}
    
    \item \textbf{Derive Amplitude $a$}:\\
    Substituting $\mu = \dfrac{\alpha^2 - \beta^2}{v^2 - u^2}$ into $v^2 = \mu a^2 - \frac{\beta^2}{\mu}$:
    \begin{equation*}
        a = \sqrt{\frac{(v^2 - u^2)(\alpha^2 v^2 - \beta^2 u^2)}{(\beta^2 - \alpha^2)^2}}
    \end{equation*}
\end{enumerate}
\end{examplebox}

\begin{examplebox}{Average Speed and Acceleration}
Show that in SHM, average speed and average acceleration magnitude over a quarter period equal $0.637$ times their maximum values.

\tcblower
\textbf{Solution:}

Let displacement be $x(t) = a \cos(\omega t)$ where $\omega = \sqrt{\mu}$.
The velocity and acceleration magnitudes at time $t$ are:
\begin{align*}
    v(t) &= |-a\omega \sin(\omega t)| = a\omega \sin(\omega t) \\
    f(t) &= |-a\omega^2 \cos(\omega t)| = a\omega^2 \cos(\omega t)
\end{align*}
The quarter period duration is $t_q = \frac{T}{4} = \frac{\pi}{2\omega}$.

\begin{enumerate}[leftmargin=1.5em]
    \item \textbf{Average Speed over Quarter Period}:\\
    \begin{align*}
        v_{\text{avg}} &= \frac{1}{t_q} \int_0^{t_q} v(t) \, dt = \frac{2\omega}{\pi} \int_0^{\frac{\pi}{2\omega}} a\omega \sin(\omega t) \, dt \\
        &= \frac{2 a \omega^2}{\pi} \left[ -\frac{\cos(\omega t)}{\omega} \right]_0^{\frac{\pi}{2\omega}} = \frac{2}{\pi} (a\omega) = \frac{2}{\pi} v_{\max} \approx 0.6366 v_{\max}
    \end{align*}
    
    \item \textbf{Average Acceleration Magnitude over Quarter Period}:\\
    \begin{align*}
        f_{\text{avg}} &= \frac{1}{t_q} \int_0^{t_q} f(t) \, dt = \frac{2\omega}{\pi} \int_0^{\frac{\pi}{2\omega}} a\omega^2 \cos(\omega t) \, dt \\
        &= \frac{2 a \omega^3}{\pi} \left[ \frac{\sin(\omega t)}{\omega} \right]_0^{\frac{\pi}{2\omega}} = \frac{2}{\pi} (a\omega^2) = \frac{2}{\pi} f_{\max} \approx 0.6366 f_{\max}
    \end{align*}
\end{enumerate}
Thus, $v_{\text{avg}} \approx 0.637 v_{\max}$ and $f_{\text{avg}} \approx 0.637 f_{\max}$.
\end{examplebox}

\begin{examplebox}{SHM with Rest Points at A(a) and B(b)}
A body moving in straight line $OAB$ with SHM has zero velocity at points $A$ and $B$ (distances $a$ and $b$ from origin $O$). It has velocity $v$ halfway between $A$ and $B$. Show that complete period is $T = \frac{\pi(b - a)}{v}$.

\tcblower
\textbf{Solution:}

\begin{enumerate}[leftmargin=1.5em]
    \item \textbf{Identify Center of Oscillation and Amplitude}:\\
    Since velocity is zero at $A(a)$ and $B(b)$, these points are the two extreme positions of oscillation.
    The center of oscillation $C$ is the midpoint of $AB$:
    \begin{equation*}
        x_C = \frac{a + b}{2}
    \end{equation*}
    The amplitude $A_0$ is half the total distance $AB$:
    \begin{equation*}
        A_0 = \frac{b - a}{2}
    \end{equation*}

    \item \textbf{Evaluate Velocity at Midpoint $C$}:\\
    The point halfway between $A$ and $B$ is the center of oscillation $C$, where speed reaches its maximum value:
    \begin{equation*}
        v = v_{\max} = A_0 \sqrt{\mu} = \left(\frac{b - a}{2}\right) \sqrt{\mu} \implies \sqrt{\mu} = \frac{2v}{b - a}
    \end{equation*}

    \item \textbf{Calculate Period $T$}:\\
    \begin{equation*}
        T = \frac{2\pi}{\sqrt{\mu}} = \frac{2\pi}{\left(\frac{2v}{b - a}\right)} = \frac{\pi(b - a)}{v}
    \end{equation*}
\end{enumerate}
\end{examplebox}

\begin{examplebox}{Return Time to a Point}
A particle performs SHM of period $T$ about center $O$. It passes through a point $P$ ($OP = b$) with velocity $v$ in the direction of $OP$. Prove that the time which elapses before it returns to $P$ is:
\begin{equation*}
    t = \frac{T}{\pi} \tan^{-1}\left( \frac{v T}{2\pi b} \right)
\end{equation*}

\tcblower
\textbf{Solution:}

\begin{enumerate}[leftmargin=1.5em]
    \item \textbf{Relation Between Amplitude and Velocity at $P$}:\\
    Let amplitude be $a$. Using $v^2 = \mu(a^2 - b^2)$:
    \begin{equation*}
        a^2 - b^2 = \frac{v^2}{\mu} = \frac{v^2}{(2\pi/T)^2} = \left(\frac{v T}{2\pi}\right)^2 \implies \sqrt{a^2 - b^2} = \frac{v T}{2\pi}
    \end{equation*}

    \item \textbf{Calculate Time $t_1$ from $P(b)$ to Extreme Position $A(a)$}:\\
    Using displacement equation $x(t) = a \cos(\sqrt{\mu} t_1)$:
    \begin{align*}
        b &= a \cos(\sqrt{\mu} t_1) \implies \cos(\sqrt{\mu} t_1) = \frac{b}{a} \\
        \tan(\sqrt{\mu} t_1) &= \frac{\sqrt{a^2 - b^2}}{b} = \frac{\left(\frac{v T}{2\pi}\right)}{b} = \frac{v T}{2\pi b}
    \end{align*}
    Taking inverse tangent:
    \begin{equation*}
        \sqrt{\mu} t_1 = \tan^{-1}\left(\frac{v T}{2\pi b}\right) \implies t_1 = \frac{1}{\sqrt{\mu}} \tan^{-1}\left(\frac{v T}{2\pi b}\right)
    \end{equation*}

    \item \textbf{Total Return Time}:\\
    The total time to travel from $P$ to extreme $A$ and back to $P$ is $t_{\text{return}} = 2 t_1$:
    \begin{equation*}
        t_{\text{return}} = 2 \left( \frac{T}{2\pi} \right) \tan^{-1}\left(\frac{v T}{2\pi b}\right) = \frac{T}{\pi} \tan^{-1}\left(\frac{v T}{2\pi b}\right)
    \end{equation*}
\end{enumerate}
\end{examplebox}

\begin{examplebox}{Period from Three Consecutive Equal Interval Measurements}
A particle moves with SHM. Its distances from center at 3 consecutive seconds are $x_1, x_2, x_3$. Prove that complete period is:
\begin{equation*}
    T = \frac{2\pi}{\cos^{-1}\left( \frac{x_1 + x_3}{2 x_2} \right)}
\end{equation*}

\tcblower
\textbf{Solution:}

\begin{enumerate}[leftmargin=1.5em]
    \item \textbf{Express Displacements at 1-Second Intervals}:\\
    Let time $t_0$ correspond to position $x_2$. Then positions at $t_0 - 1$, $t_0$, and $t_0 + 1$ are:
    \begin{align*}
        x_1 &= a \cos(\sqrt{\mu} t_0 - \sqrt{\mu}) \\
        x_2 &= a \cos(\sqrt{\mu} t_0) \\
        x_3 &= a \cos(\sqrt{\mu} t_0 + \sqrt{\mu})
    \end{align*}

    \item \textbf{Add $x_1$ and $x_3$}:\\
    Using trigonometric identity $\cos(A-B) + \cos(A+B) = 2\cos A \cos B$:
    \begin{align*}
        x_1 + x_3 &= a \left[ \cos(\sqrt{\mu} t_0 - \sqrt{\mu}) + \cos(\sqrt{\mu} t_0 + \sqrt{\mu}) \right] \\
        &= 2 a \cos(\sqrt{\mu} t_0) \cos(\sqrt{\mu}) = 2 x_2 \cos(\sqrt{\mu})
    \end{align*}

    \item \textbf{Solve for Period $T$}:\\
    \begin{equation*}
        \cos(\sqrt{\mu}) = \frac{x_1 + x_3}{2 x_2} \implies \sqrt{\mu} = \cos^{-1}\left(\frac{x_1 + x_3}{2 x_2}\right)
    \end{equation*}
    Since $T = \dfrac{2\pi}{\sqrt{\mu}}$:
    \begin{equation*}
        T = \frac{2\pi}{\cos^{-1}\left(\frac{x_1 + x_3}{2 x_2}\right)}
    \end{equation*}
\end{enumerate}
\end{examplebox}
